% Original synthetic content using the unmodified IEEEtran conference class.
\documentclass[conference]{IEEEtran}
\usepackage{amsmath,amssymb,graphicx,cite,hyperref}
\newcommand{\mean}[1]{\overline{#1}}
\newtheorem{claim}{Claim}[section]
\title{Synthetic computer science compatibility audit}
\author{\IEEEauthorblockN{Ada Audit}\IEEEauthorblockA{Department of Computing\\Synthetic University\\ada@example.invalid}\and\IEEEauthorblockN{Ben Check}\IEEEauthorblockA{Systems Institute\\Synthetic Laboratory\\ben@example.invalid}}
\begin{document}
\maketitle
\begin{abstract}
An original synthetic manuscript checks template metadata, formulas, references, figures and tables.
\end{abstract}
\begin{IEEEkeywords}
synthetic, compatibility, estimator
\end{IEEEkeywords}
\section{Synthetic estimator}\label{sec:one}
A synthetic estimator has mean $\mean{x}$. See Equation~\eqref{eq:one}, Figure~\ref{fig:one}, Table~\ref{tab:one}, and Claim~\ref{cl:one}. Our synthetic comparison cites~\cite{audit2024,check2025}.
\begin{equation}
\mean{x} = \frac{1}{n}\sum_{i=1}^n x_i.\label{eq:one}
\end{equation}
\begin{claim}[Bound]\label{cl:one}
If $0\leq x_i\leq 1$, then $0\leq\mean{x}\leq 1$.
\end{claim}

\begin{figure}
\includegraphics[width=.75\columnwidth]{signal.png}
\caption{Synthetic experiment signal.}\label{fig:one}
\end{figure}
\begin{table}
\caption{Synthetic input counts.}\label{tab:one}
\begin{tabular}{cc}
Dataset & Count\\
\hline
Alpha & 8\\
Beta & 16
\end{tabular}
\end{table}
\appendix
\section{Synthetic supplement}\label{sec:supplement}
This supplement gives no real scientific claim. Return to Section~\ref{sec:one}.
\bibliographystyle{IEEEtran}
\bibliography{refs}
\end{document}
